
\section{Contesto: Il costo della manutenzione del software e il ruolo dell'Automated Program Repair (APR)}

I difetti del software rappresentano una delle sfide principali e più costose nell'intero ciclo di vita dello sviluppo informatico, 
comportando un impiego massiccio di tempo, denaro e risorse umane \cite{le2016towards}.

Come evidenziato da diversi studi, il processo di risoluzione dei bug (bug fixing) è un'operazione notoriamente ardua, laboriosa e 
dispendiosa in termini di tempo, la quale arriva a dominare le agende degli sviluppatori e a costituire la voce di costo predominante 
nella manutenzione del software \cite{le2016towards}.

L'impatto economico derivante dalla scarsa qualità del codice è di proporzioni critiche: secondo un recente rapporto del Consortium 
for Information \& Software Quality (CISQ), citato nelle ricerche di Huang et al., nel solo 2020 gli Stati Uniti hanno sostenuto costi 
pari a circa 2,08 trilioni di dollari a causa di software di bassa qualità, di cui ben 1,56 trilioni imputabili direttamente ai fallimenti 
operativi dei sistemi (Cost of Poor Software Quality).

Tali cifre dimostrano come gli ingegneri del software debbano investire una quantità considerevole di sforzi manuali per garantire il 
corretto funzionamento dei programmi, mitigando i difetti attraverso continue e tediose attività di test e patching \cite{huang2022survey}.

Tuttavia, l'approccio tradizionale basato esclusivamente sul debugging manuale sta diventando rapidamente insostenibile. La società moderna 
dipende sempre di più dai sistemi software, il cui rapido sviluppo e la crescente popolarità su larga scala ne hanno aumentato esponenzialmente 
la complessità strutturale \cite{huang2024evolving}.
Come discusso nello studio di Dikici e Bilgin, l'avvento di nuove tecnologie e il costante aumento della scala e dell'intricatezza delle 
applicazioni moderne rendono le metodologie di ispezione umana del tutto insufficienti a soddisfare la crescente domanda di efficienza e accuratezza 
nella manutenzione.

I difetti software, che derivano per lo più da problemi storici radicati nelle fasi iniziali di progettazione, rendono la riparazione manuale
un compito sempre più prono all'errore (error-prone) e incompatibile con i cicli di sviluppo rapidi odierni \cite{huang2024evolving,dikici2025advancements}.
Di conseguenza, emerge la necessità impellente di individuare soluzioni scalabili che possano ridurre in modo drastico il carico di lavoro manuale \cite{le2016towards}.

In questo scenario critico, l'Automated Program Repair (APR) si inserisce come una soluzione trasformativa mirata a risolvere direttamente 
l'insostenibilità dei costi e dello sforzo umano appena descritti \cite{dikici2025advancements}.
Come suggerito da Huang et al., l'obiettivo principale dell'APR è quello di trasferire il pesante lavoro di manutenzione manuale verso un 
processo di riparazione automatizzato, preciso ed efficiente. 

Sfruttando avanzate tecniche computazionali, che spaziano dalla tradizionale ricerca euristica (search-based) ai metodi basati su vincoli e 
template, fino ad arrivare ai più recenti approcci guidati dai dati tramite apprendimento automatico e profondo (deep learning), i sistemi APR 
sono specificamente progettati per analizzare il codice e rilevare le anomalie in autonomia \cite{huang2024evolving,dikici2025advancements}.

Il loro scopo ultimo è quello di generare e implementare patch valide con un intervento umano minimo o del tutto assente \cite{dikici2025advancements}.

Automatizzando il tedioso processo del bug fixing, le tecnologie APR si posizionano come una risorsa cruciale: non solo accelerano i cicli di 
sviluppo riducendo sensibilmente le tempistiche e i costi necessari per risolvere i difetti, ma contribuiscono in modo determinante alla 
creazione di sistemi software più robusti, affidabili e di alta qualità \cite{dikici2025advancements,le2016towards}.